\subsection{World Models: The Deep Revival}
\label{section:fc_ha_schmidhuber}

\citet{HaSchmidhuber2018} is the first paper to make latent-space dreaming work at scale on pixel inputs. An agent learns a visual encoder, a recurrent forecaster of the encoder's output, and a controller, with the controller trained against rollouts of the forecaster rather than against real environment frames. The abstract decomposition is a triple of objects with distinct losses, a vision encoder trained against reconstruction, a memory model trained against next-latent likelihood, and a controller trained against imagined returns; this decomposition is the architectural pattern that subsequent deep model-based methods inherit and modify, with the Dreamer line collapsing vision and memory into one jointly trained object and the value-aware line of \S\ref{section:fc_value_aware} replacing the memory loss outright. On the CarRacing benchmark the resulting controller is the first to cross the conventional solution threshold; on a VizDoom task a controller trained entirely inside the learned forecaster transfers zero-shot to the real environment. The paper is an explicit homage to \citet{Schmidhuber1990}, the controller-model architecture introduced in \S\ref{section:fc_origins_schmidhuber}, and the line of work it revives is twenty-eight years old. What is new is the toolkit. Convolutional variational autoencoders, mixture-density recurrent networks, and evolution strategies, taken together, make the older architecture empirically tractable at a scale the original could not reach.

Nine recurring terms anchor the rest of the chapter. A \emph{world model} is an estimated model of next states and rewards. A \emph{planning update} is a policy or value update performed using the learned model rather than the environment. A \emph{simulated rollout} is a trajectory generated by the learned model rather than by real interaction. A \emph{latent state} is an internal compressed state used by a neural model, not an observed economic state. A \emph{model loss} is the criterion used to train the world model. A \emph{value-aware loss} is a model loss that cares about preserving values rather than predicting every detail of the next observation. A \emph{short rollout} is a deliberately short simulated trajectory from a real visited state. A \emph{bootstrap} is a value estimate that stands in for rewards beyond the simulated horizon. A \emph{model error} is the gap between the learned model and the true environment, judged by its effect on planning rather than by its statistical fit.

\subsubsection{Three components}
\label{section:fc_ha_schmidhuber_components}

The vision component $V$ is a convolutional variational autoencoder.\footnote{A variational autoencoder compresses high-dimensional observations into a low-dimensional latent vector by training an encoder and a decoder jointly.} The encoder maps an observation $x_t \in \mathbb{R}^{64 \times 64 \times 3}$ to a latent vector $z_t \sim \mathcal{N}(\mu(x_t), \sigma(x_t)^2 I)$ with $z_t \in \mathbb{R}^{N_z}$, where $N_z = 32$ for CarRacing and $N_z = 64$ for VizDoom. The decoder is trained jointly with the encoder against the standard ELBO objective,\footnote{Evidence lower bound; the standard variational autoencoder training objective that combines a reconstruction term with a Kullback-Leibler penalty between the encoder's latent and a fixed prior. The Kullback-Leibler term measures the discrepancy between two distributions.} combining a per-pixel reconstruction term with a Kullback-Leibler penalty on the latent. The vision component is trained on observation reconstruction alone, independently of the controller, on a buffer of frames collected by a random policy. The decoder is discarded once training is complete; downstream components see only the encoder's latents. The architectural commitment is that all subsequent processing happens in a low-dimensional vector space rather than in pixel space, which is what makes the rest of the pipeline cheap enough to train and to roll out.

The memory component $M$ is a mixture-density recurrent neural network.\footnote{A mixture-density recurrent network is a recurrent network whose output parameterizes a Gaussian mixture over the next latent rather than a single mean; this matters when the next state is genuinely multimodal.} Its body is an LSTM with hidden state $h_t$;\footnote{Long short-term memory; a gated recurrent cell that propagates information across time and is the dominant pre-Transformer architecture for sequential data.} its head outputs the parameters of a Gaussian mixture over the next latent, conditional on the current action, the current latent, and the recurrent state, modeling $P(z_{t+1} \mid a_t, z_t, h_t)$ as a mixture of five diagonal Gaussians. A temperature parameter on the mixture controls how aggressively the forecaster samples from its tails, and is exposed as a tuning lever for controller training. The memory component is trained on next-latent likelihood alone, again independently of the controller, using the same buffer of rollouts under a random policy that trained $V$. The mixture form matters in environments with discrete event structure such as enemy fireballs in VizDoom, where the next latent has multiple plausible values that a unimodal Gaussian forecaster would average over and so render unusable for planning.

The controller $C$ is a single linear layer. It outputs an action $a_t = W_c [z_t, h_t] + b_c$ from the concatenation of the visual latent and the recurrent state. On CarRacing the controller has roughly $870$ parameters, compared to about $5$ million in $V$ and $M$ together. It is trained by covariance matrix adaptation evolution strategy,\footnote{A gradient-free black-box optimizer that maintains a Gaussian over candidate parameter vectors and updates the Gaussian by re-fitting to the best-scoring samples each generation.} evaluating candidate parameter vectors against rollouts drawn from $M$ rather than against real environment frames. Evolution strategies sidestep the need for differentiability through $M$ at the cost of higher sample complexity per parameter, which the linear policy keeps manageable. The controller is the only component whose loss depends on reward, and it is the only component whose training requires interaction, in the form of imagined trajectories through the memory network.

\subsubsection{Empirical headline}
\label{section:fc_ha_schmidhuber_empirics}

On CarRacing-v0 the full agent scores $906 \pm 21$ over $100$ random tracks. The benchmark requires an average above $900$ to count as solved, and this is the first reported solution. On the VizDoom DoomTakeCover task the controller is trained inside the memory network, with the visual encoder and recurrent forecaster fitted on a buffer of random rollouts and the controller never exposed to real frames during policy search. Transferred zero-shot to the real environment, the controller scores $1092 \pm 556$ against a solution threshold of $750$. The variance is large but the mean clears the threshold, and the policy is constructed without any direct contact with the real environment after the data collection phase.

\subsubsection{Architectural reading}
\label{section:fc_ha_schmidhuber_reading}

The modular training schedule, with $V$ fitted on reconstruction, $M$ fitted on next-latent likelihood, and $C$ fitted on dream rollouts, is the practical workaround for joint end-to-end training, which the paper notes is hard at this scale. Each component has its own loss, its own optimizer, and its own data pipeline, and the three are composed in sequence rather than co-trained. The cost of the workaround is that the representation learned by $V$ is agnostic to what the controller will use it for, and the forecaster learned by $M$ is agnostic to what the controller will reward. The line of work that eventually unifies $V$ and $M$ into a single jointly-trained recurrent state-space model is the PlaNet and Dreamer family, treated in \S\ref{section:fc_rssm}.

The combination of $V$ and $M$ does the representation work; the controller is tiny by comparison and carries the parameter budget of a linear policy. This anticipates the modern observation that a strong world model lets a small actor produce competitive policies, a point pressed further by Dreamer in \S\ref{section:fc_rssm} and by TD-MPC2 in \S\ref{section:fc_tdmpc2}, where the latent space and the value function carry most of the inductive load and the policy network is relatively light. The training-in-dream-then-transfer protocol on VizDoom is the strongest demonstration to date that a learned forecaster can stand in for the environment during policy search, with the agent's only contact with the real environment being the random data collection used to fit $V$ and $M$ in the first place.
